\section{Discussion and Conclusion}
\label{sec:conclusion}

A common way of training VLMs for corruption robustness quietly makes them worse at
admitting when they cannot see; standard accuracy benchmarks miss the regression; and it
is recoverable. The failure is behavioral, not representational (RQ1): answerability
stays linearly decodable after fine-tuning. The two dominant recipes fail for opposite
reasons (RQ2): RKL on the reasoning side, GRPO on the answer side via a reward-induced
shortcut. And because the knowledge survives, it can be put back to work (RQ3): LATCH reads it
before generation and restores correct refusals on every fine-tuned model, at a cost of at
most $1.3$ points of corrupted accuracy, while the corruption gains still transfer
zero-shot.

That fine-tuning for one capability silently erodes another, in a direction invisible to
the metric being optimized, echoes what is seen when fine-tuning erodes safety
alignment~\cite{qi2024finetuning}, and is a pattern we expect to recur wherever a training
signal is defined only over answerable inputs. The practical message is simple but
consequential: wherever a confident wrong answer is dangerous, corruption robustness and
calibrated abstention must be measured, and optimized, together.

\paragraph{Limitations.}
We study three 4B backbones and four benchmarks, with one training seed per configuration;
whether the two failure signatures generalize to other scales, reward designs and
corruption regimes is open. Refusals are detected by keyword matching on the final answer,
which can miss paraphrased refusals and count hedged answers as refusals. The answerability
probe and the reasoning transplant were run on Qwen3-VL-4B only. LATCH reads a single
layer with a fixed cut-off and judges the input rather than the answer, so it can override
a correct reading of a blurry but legible image; adaptive, multi-layer cut-offs are natural
extensions.
